\section{Background}

\subsection{Reinforcement Learning from Images} 
We formulate image-based control as an infinite-horizon Markov Decision Process (MDP)~\citep{bellman1957mdp}. Generally, in such a setting, an image rendering of the system is not sufficient to perfectly describe the system's underlying state. To this end and per common practice~\citep{mnih2013dqn}, we approximate the current state of the system by stacking three consecutive prior observations. With this in mind, such MDP can be described as a tuple $(\gX, \gA, P, R, \gamma, d_0)$, where $\gX$ is the state space (a three-stack of image observations), $\gA$ is the action space, $P: \gX \times \gA \to \Delta(\gX)$ is the transition function\footnote{Here, $\Delta(\gX)$ denotes a distribution over the state space $\gX$.} that defines a probability distribution over the next state given the current state and action, $R: \gX \times \gA \to [0,1]$ is the reward function, $\gamma \in [0, 1)$ is a discount factor, and $d_0 \in \Delta(\gX)$ is the distribution of the initial state $\vx_0$. The goal is to find a policy $\pi: \gX \to \Delta(\gA)$ that maximizes the expected discounted sum of rewards $\E_\pi[\sum_{t=0}^\infty \gamma^t r_t]$, where  $\vx_0 \sim d_0$, and $\forall t$ we have $\va_{t} \sim \pi(\cdot|\vx_{t})$, $\vx_{t+1} \sim P(\cdot| \vx_{t}, \va_{t})$, and $r_t = R(\vx_{t}, \va_{t})$.



\subsection{Deep Deterministic Policy Gradient}
Deep Deterministic Policy Gradient (DDPG)~\citep{lillicrap2015ddpg} is an actor-critic algorithm for continuous control that concurrently learns a Q-function $Q_\theta$ and a deterministic policy $\pi_\phi$. For this, DDPG uses Q-learning~\citep{watkins1992qlearning} to learn $Q_\theta$ by minimizing the one-step Bellman residual 
    $J_\theta(\gD) = \E_{\substack{( \vx_t,\va_t, r_t, \vx_{t+1}) \sim \gD }}[(Q_\theta(\vx_t, \va_t) - r_t - \gamma Q_{\bar{\theta}}(\vx_{t+1},\pi_\phi(\vx_{t+1} ) )^2]$.
The policy $\pi_\phi$ is learned by employing Deterministic Policy Gradient (DPG)~\citep{silver14dpg} and maximizing $J_\phi(\gD) = \E_{\vx_t \sim \gD} [Q_\theta(\vx_t, \pi_\phi(\vx_t))]$, so $\pi_\phi(\vx_t)$ approximates $\mathrm{argmax}_\va Q_\theta(\vx_t, \va)$.  Here, $\gD$ is a replay buffer of environment transitions and $\bar{\theta}$ is an exponential moving average of the weights. DDPG is amenable to incorporate $n$-step returns~\citep{watkins1989phd,peng1996nstepq} when estimating TD error beyond a single step~\citep{barth-maron2018d4pg}. In practice, $n$-step returns allow for faster reward propagation and has been previously used in policy gradient and Q-learning methods~\citep{mnih2016asynchronous,barth-maron2018d4pg,hessel2017rainbow}.

\subsection{Data Augmentation in Reinforcement Learning}
Recently, it has been shown that data augmentation techniques, commonplace in Computer Vision, are also important for achieving the state-of-the-art performance in image-based RL~\citep{yarats2021image,laskin2020reinforcement}. For example, the state-of-the-art algorithm for visual RL, DrQ~\citep{yarats2021image} builds on top of Soft Actor-Critic~\citep{haarnoja2018sac}, a model-free actor-critic algorithm, by adding a convolutional encoder and data augmentation in the form of random shifts. The use of such data augmentations now forms an essential component of several recent visual RL algorithms~\citep{srinivas2020curl,raileanu2020automatic,yarats2021proto,stooke2020decoupling,hansen2021softda,schwarzer2020spr}. 

